\section*{第 \ref{ch:g10:lines} 章——直线方程与线性方程组}

\begin{solution}{exo:g10:lines:1}
$y = 2x - 3$：\omterm{def:g10:reffunc:affine}{斜率} $2$，截距 $-3$。
$y = -\frac13 x + 1$：\omterm{def:g10:reffunc:affine}{斜率} $-\frac13$，截距 $1$。
$y = 4$：\omterm{def:g10:reffunc:affine}{斜率} $0$，截距 $4$（水平直线）。
$x = -2$：竖直直线，没有\omterm{def:g10:reffunc:affine}{斜率}，也没有\omterm{thm:g10:lines:reduced}{斜截式方程}。
\end{solution}

\begin{solution}{exo:g10:lines:2}
$3 \times 2 + 1 = 7$：是，点 $A$ 在直线上。
$3 \times (-1) + 1 = -2 \neq -1$：点 $B$ 不在直线上。
$3 \times 4 + 1 = 13$：是，点 $C$ 在直线上。
\end{solution}

\begin{solution}{exo:g10:lines:3}
(a) $m = \frac{8 - 2}{3 - 0} = 2$，且 $p = y_A = 2$（点 $A$ 在
$y$ 轴上）：$y = 2x + 2$。

(b) $m = \frac{-4 - 5}{2 - (-1)} = \frac{-9}{3} = -3$，
由 $5 = -3 \times (-1) + p$ 得 $p = 2$：$y = -3x + 2$。

(c) $x_A = x_B = 4$：竖直直线 $x = 4$。
\end{solution}

\begin{solution}{exo:g10:lines:4}
$y = \frac{6x+2}{2} = 3x + 1$。\omterm{def:g10:reffunc:affine}{斜率}为 $3$ 的直线有 $y = 3x - 1$，
$y = 3x + 5$ 和 $y = 3x + 1$，这三条直线两两平行；
三者两两不同（截距各不相同），且都不与
$y = -3x + 2$ 平行。
\end{solution}

\begin{solution}{exo:g10:lines:5}
第一个\omterm{def:g10:algebra:equation}{方程}组：将 $y = 2x - 1$ 代入 $3x + y = 9$，
得 $3x + 2x - 1 = 9$，即 $5x = 10$，$x = 2$，进而 $y = 3$。解为
$(2, 3)$。

第二个\omterm{def:g10:algebra:equation}{方程}组：由 $x - y = 4$ 得 $x = y + 4$。代入得
$2(y + 4) + 3y = 3$，即 $5y = -5$，$y = -1$，进而 $x = 3$。解为
$(3, -1)$。
\end{solution}

\begin{solution}{exo:g10:lines:6}
第一个\omterm{def:g10:algebra:equation}{方程}组：两式相加消去 $y$，得 $2x = 14$，$x = 7$；
再由 $7 + y = 10$ 得 $y = 3$。解为 $(7, 3)$。

第二个\omterm{def:g10:algebra:equation}{方程}组：第一个\omterm{def:g10:algebra:equation}{方程}两边乘以 $5$，第二个\omterm{def:g10:algebra:equation}{方程}两边乘以 $-2$：
\[
\begin{cases}
15x + 10y = 60 \\
-4x - 10y = -16
\end{cases}
\]
相加得 $11x = 44$，故 $x = 4$；再由 $3 \times 4 + 2y = 12$ 得
$y = 0$。解为 $(4, 0)$。
\end{solution}

\begin{solution}{exo:g10:lines:7}
分别计算各情形下的 $ab' - a'b$。

第一种情形：$2 \times 2 - (-4) \times (-1) = 4 - 4 = 0$，且第二个\omterm{def:g10:algebra:equation}{方程}
是第一个\omterm{def:g10:algebra:equation}{方程}的 $-2$ 倍，两\omterm{def:g10:algebra:equation}{方程}表示同一条直线，
有无穷多解。

第二种情形：同样有 $ab' - a'b = 0$，但右端常数不满足 $-2$ 倍关系，
即 $1 \neq -6$，这是两条不同的平行直线，无解。

第三种情形：$2 \times 1 - 1 \times (-1) = 3 \neq 0$，恰有一组解。
\end{solution}

\begin{solution}{exo:g10:lines:8}
设 $x$ 为儿童票的张数，$y$ 为成人票的张数，
则得\omterm{def:g10:algebra:equation}{方程}组：
\[
\begin{cases}
x + y = 200 \\
5x + 9y = 1352 .
\end{cases}
\]
由第一个\omterm{def:g10:algebra:equation}{方程}得 $x = 200 - y$，代入
得 $5(200 - y) + 9y = 1352$，即 $1000 + 4y = 1352$，$y = 88$，进而
$x = 112$。故儿童票 $112$ 张、成人票 $88$ 张。验算：
$5 \times 112 + 9 \times 88 = 560 + 792 = 1352$。
\end{solution}

\begin{solution}{exo:g10:lines:9}
\emph{1.} 平行即\omterm{def:g10:reffunc:affine}{斜率}同为 $2$，故可设 $y = 2x + p$，
其中 $4 = 2 \times 1 + p$，得 $p = 2$，即 $d'\colon y = 2x + 2$。

\emph{2.} 在 $x$ 轴上，$y = 0$，由 $2x + 2 = 0$ 得 $x = -1$。交点
为 $(-1, 0)$。
\end{solution}

\begin{solution}{exo:g10:lines:10}
\emph{1.} $[BC]$ 的\omterm{prop:g10:coordgeom:midpoint}{中点} $A' = \left(\frac{6+2}{2},
\frac{2+4}{2}\right) = (4, 3)$。过 $A(0,0)$ 的中线\omterm{def:g10:reffunc:affine}{斜率}为
$\frac{3 - 0}{4 - 0} = \frac34$，\omterm{def:g10:algebra:equation}{方程}为 $y = \frac34 x$。

$[AC]$ 的\omterm{prop:g10:coordgeom:midpoint}{中点}为 $B' = (1, 2)$。过 $B(6, 2)$ 的中线\omterm{def:g10:reffunc:affine}{斜率}为
$\frac{2 - 2}{1 - 6} = 0$，即水平直线 $y = 2$。

\emph{2.} 交点满足 $\frac34 x = 2$，得 $x = \frac83$，故
$G\left(\frac83, 2\right)$。第三条中线连接 $C(2,4)$ 与
$[AB]$ 的\omterm{prop:g10:coordgeom:midpoint}{中点} $C' = (3, 1)$，其\omterm{def:g10:reffunc:affine}{斜率}为
$\frac{1 - 4}{3 - 2} = -3$，\omterm{def:g10:algebra:equation}{方程}为 $y = -3x + 10$。当 $x = \frac83$ 时，
$y = -8 + 10 = 2$，故 $G$ 在该直线上。事实上
$G = \left(\frac{0 + 6 + 2}{3}, \frac{0 + 2 + 4}{3}\right)$：重心
的坐标即三个顶点坐标的平均数。
\end{solution}

\begin{solution}{pb:g10:lines:1}
\textbf{1.} \omterm{def:g10:reffunc:affine}{斜率} $\frac{8 - 2}{3 - 1} = 3$；过 $(1, 2)$ 的直线为：
$y = 3x - 1$。

\textbf{2.} $y = 3x - 1$ 与 $y = 3x + 4$：\omterm{def:g10:reffunc:affine}{斜率}相同，平行
（且不重合）。联立 $y = 3x - 1$ 与 $y = -x + 7$：
$3x - 1 = -x + 7$ 得 $x = 2$，$y = 5$，交点为 $(2, 5)$。

\textbf{3.} 两点都满足 $x = 2$，故直线竖直，
\omterm{def:g10:algebra:equation}{方程}为 $x = 2$。斜截式 $y = mx + p$ 所回答的是“一个 $y$
对应每个 $x$”的问题，而竖直直线不满足此条件：其\omterm{def:g10:reffunc:affine}{斜率}将为
无穷大。

\textbf{4.} $3 \times 4 - 1 = 11$：是的，$(4, 11)$ 在该
直线上。$3 \times (-1) - 1 = -4 \neq -5$：不在。

\textbf{5.} $y = -2x + 3$，与坐标轴交于 $(0, 3)$ 与
$\left(\frac32, 0\right)$。三角形面积：
$\frac12 \times \frac32 \times 3 = \frac94 = 2.25$。

\textbf{6.} 代入得 $3x + 2x - 3 = 7$，$x = 2$，
$y = 1$，解为 $(2, 1)$。

\textbf{7.} 两式相加得 $7x = 7$，$x = 1$；再由 $3y = 6$ 得 $y = 2$，
解为 $(1, 2)$。

\textbf{8.} 第一个\omterm{def:g10:algebra:equation}{方程}组：第二个\omterm{def:g10:algebra:equation}{方程}是第一个\omterm{def:g10:algebra:equation}{方程}的两倍，
即同一条直线被计了两次，解有无穷多个
（即 $2x - y = 3$ 上的全部点）。第二个\omterm{def:g10:algebra:equation}{方程}组：左边部分成比例，
右边不成比例（$2 \times 3 = 6 \neq 1$）：两条
平行线，无解。三分情形：\omterm{def:g10:reffunc:affine}{斜率}不同
$\to$ 一个交点；\omterm{def:g10:reffunc:affine}{斜率}相等、截距不同 $\to$
平行，无解；全部相同 $\to$ 同一直线，解
有无穷多个。

\textbf{9.} 两直线的方向\omterm{def:g10:vectors:vector}{向量}分别为 $(-b, a)$ 与
$(-d, c)$（或\omterm{def:g10:reffunc:affine}{斜率} $-\frac ab$ 与 $-\frac cd$），两\omterm{def:g10:vectors:vector}{向量}\omterm{def:g10:vectors:collinear}{共线}
当且仅当 $ad - bc = 0$ 时成立——这正是
\cref{pb:g10:vectors:1}中的行列式换了个角色。代入检验，第 7 题：
$2 \times (-3) - 3 \times 5 = -21 \neq 0$：唯一解。
第 8 题：两个\omterm{def:g10:algebra:equation}{方程}组均满足 $2 \times (-2) - (-1) \times 4 = 0$：
平行或重合，常数项区分了这两种
情形。

\textbf{10.} $ad - bc = 6 - 2m$，当 $m = 3$ 时为零。此时
\omterm{def:g10:algebra:equation}{方程}组为 $x + 3y = 3$ 与 $2x + 6y = 7$：将第一式加倍得
$2x + 6y = 6 \neq 7$，两直线平行，\emph{无解}。当
$m \neq 3$ 时，消元得唯一解
$y = \frac{1}{6 - 2m}$，再得 $x = 3 - \frac{m}{6 - 2m}$。

\textbf{11.} $p + r = 35$，$2p + 4r = 94$：两边折半得
$p + 2r = 47$，于是 $r = 12$，$p = 23$：二十三只雉，
十二只兔——正是\emph{《九章算术》}'自己的答案。

\textbf{12.} $x + y = 30$ 与 $0.6x + 0.2y = 0.35 \times 30 =
10.5$，第二式减去 $0.2 \times$ 第一式，得 $0.4x = 4.5$：
$x = 11.25$ 升糖浆，$y = 18.75$ 升饮料。

\textbf{13.} 数字 $a$、$b$：$a + b = 12$ 且
$(10a + b) - (10b + a) = 9(a - b) = 18$，故 $a - b = 2$：
$a = 7$，$b = 5$，此数为 $75$。

\textbf{14.} 两式相加得 $5c + 5r = 18$，故 $c + r = 3.60$——
即一杯咖啡加一个牛角包。相减得 $c - r = 0.20$。于是
$c = 1.90$，$r = 1.70$ 欧元：利用对称的捷径，
无需任何代入即可求解。

\textbf{15.} 将第一条说法加倍，得 $4$ 个苹果 $+ 2$
个梨 $= 10$ 欧元，但摊位收费为 $9$ 欧元：此\omterm{def:g10:algebra:equation}{方程}组
矛盾——两直线平行，解集为空。结论：
两条价格说法不可能同时正确；要么其中藏着折扣，
要么是算术有误。

\textbf{16.} 在原点处：$0 > 3 \times 0 - 1 = -1$，成立，
故 $(0,0)$ 位于区域 $y > 3x - 1$ 内（直线上方）。
任取一点代入检验：上方或下方取决于
不等式是否成立。

\textbf{17.} 顶点为：$(0, 0)$；$(4, 0)$（坐标轴与
$x + y = 4$ 的交点）；$(0, 1)$（坐标轴与 $y = 2x + 1$ 的交点）；以及
$x + y = 4$ 与 $y = 2x + 1$ 的交点：$x = 1$，$y = 3$，即 $(1, 3)$。

\textbf{18.} $P(0,0) = 0$；$P(4, 0) = 12$；$P(0, 1) = 2$；
$P(1, 3) = 9$。\omterm{def:g10:functions:extrema}{最大值}在顶点 $(4, 0)$ 处取得：方案
$x = 4$，$y = 0$ 获得 $12$——顶点决定最优，正如
平行线族 $3x + 2y = c$ 扫过时所显示的那样。

\textbf{19.} 由 $1.5x + 3 = 2x$ 得 $x = 6$，$y = 12$：
在六千米处两辆出租车收费均为十二欧元——
这正是初中分册“两支温度计”周末问题中盈亏平衡的交点，如今化为\omterm{def:g10:algebra:equation}{方程}组的唯一
解。

\textbf{20.} 二元一次\omterm{def:g10:algebra:equation}{方程}的解构成一条直线；\omterm{def:g10:lines:system}{二元一次方程组}
则叠加两条直线，$ad - bc$ 一眼即可说明
它们是否恰有一个交点（非零时），或落入
平行或重合的情形（为零）。若改为叠加不等式，
解便构成一个多边形区域，其顶点
承载任意线性最优值。线性代数将
行列式推广到任意多个未知数；线性规划
则把这些顶点付诸实用。
\end{solution}
